% !TEX root = CSL-main.tex
\section{Transition Graphs of Cyclic Proofs}
\label{sec:transition-graph-proofs}

Given a $\CLKIDomega$ or $\CLJIDomega$ pre-proof $\Pi$, we show how to construct the induced transition graph with bi-partite relations. Our approach follows Ikebuchi's work \cite{Mirai}, where she formally stated the equivalence between the GTC and SCT, by extracting the size-change graphs of a cyclic proof. In our case, we extract a transition graph with bi-partite relations of a cyclic proof, which in principle corresponds to the same process.

\begin{definition}[Transition Graph of a Pre-Proof]
\label{def:preProofTrans}

Let $\Pi = ((V, s, r, p), \mathcal{F})$ be a $\CLKIDomega$ pre-proof with a graph $G_{\Pi} = (V',s,r,p')$. We note $\mathcal{G}_{\Pi} = (V', \mathcal{R}, E)$ as the transition graph over bi-partite relations induced by $\Pi$, where:
\begin{itemize}

     \item $\mathcal{R} \subseteq X \times X$ where $X = \{ (P\mathbf{u},j) \ | \ \text{$P$ is inductive, $\exists v \in V'$ s.t. $s(v) = \Gamma \vdash \Delta$ and $(P\mathbf{u},j) \in \Gamma$} \}$

     \item Let $v,v' \in V'$ such that both vertices share and edge, i.e. there exists a $j$ such that $p'_j(v) = v'$. Moreover, let $s(v) = \Gamma \vdash \Delta$ (resp. $s(v') = \Gamma' \vdash \Delta'$) where $\Psi \subseteq \Gamma$ (resp. $\Theta \subseteq \Gamma'$)
     contains only the ocurrences of $\Gamma$ (resp. $\Gamma'$), of the form $(P\mathbf{u}, i)$, where $P$ is an inductive predicate. We define the bi-partite relation $(S,P)_{v,v'}$ as follows:
        \begin{enumerate}[i)]
            \item The domain and codomain of $|(S,P)_{v,v'}|$ are $\Psi$ and $\Theta$ respectively.
            

            \item If $r(v)$ is not a case split rule, then for all $(P\mathbf{u},i)$ in $\Psi$ and $(Q\mathbf{t},j)$ in $\Theta$, $(P\mathbf{u},i) S (Q\mathbf{t},j)$ if and only if $(P\mathbf{u},i) \rightarrow_{r(v)} (Q\mathbf{t},j)$.

            \item If $r(v)$ is a case split rule, then for all $(Q\mathbf{t},j)$ in $\Theta$, $(P\mathbf{u},i) P (Q\mathbf{t},j)$ if and only if $(P\mathbf{u},i) \rightarrow_{r(v)} (Q\mathbf{t},j)$ and $(P\mathbf{u},i)$ is the active occurrence in $\Gamma$. Otherwise, if $(P\mathbf{u},i)$ is not the active occurrence, then $(P\mathbf{u},i) S (Q\mathbf{t},j)$.

        \end{enumerate}

    \item $\mathcal{R} = \{(S,P)_{v,v'} \ | \ \text{for all $v,v' \in V$ s.t. $p'_j(v) = v'$ for some $j$}  \}$

    \item $E = \{\edge{(S,P)_{v,v'}}{v}{v'} \ | \ \text{for all $v,v' \in V$ s.t. $p'_j(v) = v'$ for some $j$}  \}$
 \end{itemize} 

This process can be also applied to $\CLJIDomega$ pre-proofs, using sequents of the form $\Gamma \vdash \phi$.

\end{definition}

Note that this construction defines an isomorphism between graphs. Specifically, $G_{\Pi}$ is isomorphic to $\mathcal{G}_{\Pi}$, in the sense that for all $v,v' \in V'$ then $p'_j(v) = v'$ in $G_{\Pi}$ if and only if $\edge{(S,P)_{v,v'}}{v}{v'}$ in $\mathcal{G}_{\Pi}$.


% !TEX root = CSL-main.tex
\begin{figure}[]
\centering
\resizebox{0.75\textwidth}{!}{%
\begin{circuitikz}
\tikzstyle{every node}=[font=\LARGE]
\node [font=\LARGE] at (2,12.25) {$v_0$};
\node [font=\LARGE] at (5.75,12.25) {$v_1$};
\node [font=\LARGE] at (9.5,13.75) {$v_2$};
\node [font=\LARGE] at (9.5,11) {$v_3$};
\node [font=\LARGE] at (13.25,11) {$v_4$};
\node [font=\LARGE] at (17,11) {$v_5$};
\node [font=\LARGE] at (20.75,11) {$v_6$};
\draw [->, >=Stealth] (3,12.25) -- (5,12.25);
\draw [->, >=Stealth] (6.25,12.25) -- (8.75,13.5);
\draw [->, >=Stealth] (6.25,12) -- (8.75,11);
\draw [->, >=Stealth] (10.25,11) -- (12.5,11);
\draw [->, >=Stealth] (14,11) -- (16.25,11);
\draw [->, >=Stealth] (17.75,11) -- (20,11);
\node [font=\large] at (4,12.75) {$(\{ \textcolor{blue}{(Nx, Nx)}\}, \textcolor{red}{\emptyset})$};
\node [font=\large] at (7.25,13.5) {$( \textcolor{blue}{\emptyset}, \textcolor{red}{\emptyset})$};
\node [font=\large] at (7.25,10.75) {$( \textcolor{blue}{\emptyset}, \{\textcolor{red}{(Nx,Ny)}\})$};
\node [font=\large] at (11.25,10.5) {$(\{ \textcolor{blue}{(Ny, Ny)}\}, \textcolor{red}{\emptyset})$};
\node [font=\large] at (15,10.5) {$(\{ \textcolor{blue}{(Ny, Ny)}\}, \textcolor{red}{\emptyset})$};
\node [font=\large] at (18.75,10.5) {$(\{ \textcolor{blue}{(Ny, Ny)}\}, \textcolor{red}{\emptyset})$};
\draw [->, >=Stealth] (20.75,10.5) .. controls (16.25,8.5) and (6.25,7.5) .. (2.25,11.5) ;
\node [font=\large] at (11,8) {$(\{ \textcolor{blue}{(Ny, Nx)}\}, \textcolor{red}{\emptyset})$};
\end{circuitikz}
}
\caption{Transition graph of the cyclic pre-proof in Figure \ref{exampleCircProof}. For each relation $(S,P)_{v,v'}$, the relation $S$ is represented by blue pairs, while pairs in $P$ are colored in red.}

\label{exampleTransGraph}
\end{figure}

\begin{example}
    Using Definition \ref{def:preProofTrans} over the cyclic pre-proof in Figure \ref{exampleCircProof}, we obtain the transition graph in Figure \ref{exampleTransGraph}. 
\end{example}

Using the construction specified in Definition \ref{def:preProofTrans}, and the Definitions \ref{def:TracePath} and \ref{def:traceCyclic} of path and traces for transitions graphs and cyclic pre-proofs, it can be derived that:


\begin{proposition}
Consider a circular pre-proof $\Pi$ with a graph $G_{\Pi}$ and induced transition graph $\mathcal{G}_{\Pi}$. Moreover, let $\pi$ be a sequence of vertices defining a path in $G_{\Pi}$, and $\pi'$ corresponding to the same path in $\mathcal{G}_{\Pi}$, but defined as in Definition \ref{def:TracePath} (we can do this, since both graphs are isomorphic). Then, $\uptau$ is a trace over $\pi$ if and only if $\uptau \models \pi'$.
\end{proposition}

Then, using the previous proposition and the Definitions \ref{def:GTC} and \ref{def:SCT}, the equivalence between GTC and SCT follows:

\begin{proposition}
Let $\Pi$ be a $\CLKIDomega$ (resp. $\CLJIDomega$) pre-proof. Then $\Pi$ is a $\CLKIDomega$ (resp. $\CLJIDomega$) proof if and only if $\mathcal{G}_{\Pi}$ satisfies the size-change termination condition.
\end{proposition}